\section{Cyclotomic-unit matching and a scalar shift resultant}

Put $u_n=s_n/s_1$.  For valid trace classes define
\begin{equation}\label{eq:WV}
W_r(b)=u_{(b+r)/2}u_{(b-r)/2},\qquad
V_r(a)=u_ru_{(a+1)/2}u_{(a-1)/2}.
\end{equation}
A direct difference factorization gives
\begin{equation}\label{eq:unitmatch}
\gamma_r(a,b)-\beta_r=s_1^2\bigl(W_r(b)-V_r(a)\bigr).
\end{equation}
Since $s_1$ is a unit away from $K$, every split row intersection is exactly a
congruence $W_r(b)\equiv V_r(a)\pmod{\mathfrak q}$ between explicit products
of cyclotomic units.

Let
\[
P_1(T)=\frac{H_K(T)}{T-t_1},\qquad
R_r(U)=\frac{H_K(U)}{U-t_r}
\]
and define the scalar shift resultant
\begin{equation}\label{eq:shiftres}
\mathcal R_{K,r}(Z)
 =\Res_T\!\left(P_1(T),R_r(\alpha_rT+\beta_r+Z)\right)
 =c_0+c_1Z+c_2Z^2+\cdots.
\end{equation}

\begin{theorem}[Scalar support identity]\label{thm:scalar-support}
For every split embedding prime $\mathfrak q$,
\begin{equation}\label{eq:ordshift}
\ord_{Z=0}(\mathcal R_{K,r}(Z)\bmod\mathfrak q)
 =I_{1,r}(\mathfrak q).
\end{equation}
Consequently
\[
I_{1,r}(\mathfrak q)\ge3
\iff
c_0\equiv c_1\equiv c_2\equiv0\pmod{\mathfrak q}.
\]
\end{theorem}

\begin{proof}
Over the split residue field the two valid trace polynomials factor linearly,
so the product formula gives
\[
\mathcal R_{K,r}(Z)
 =\prod_{[a]\ne[1]}\prod_{[b]\ne[r]}
   (Z+\alpha_rt_a+\beta_r-t_b).
\]
Each factor is linear in $Z$ and vanishes at zero exactly for one row
intersection.
\end{proof}

The ideal $(c_0,c_1,c_2)$ contracts to an integer
$\Omega_{K,r}^{(3)}$ with the same split rational-prime order-three support as
$\Lambda_{K,r}^{(3)}$.  Equality of the full integers is not asserted.
Moreover, the nonzero characteristic-zero constant coefficient gives the
per-ratio deterministic budget, where
$I_r(p):=\max_{\mathfrak q\mid p}I_{1,r}(\mathfrak q)$ over the embedding
primes $\mathfrak q$ of the split prime $p$:
\begin{equation}\label{eq:perratio}
\sum_{\substack{p\equiv1\,(K)\\I_r(p)\ge3}}\log p=O(K^3).
\end{equation}
Indeed, $c_0\ne0$ in characteristic zero: by
Theorem~\ref{thm:intercept} the distinguished intercept is isolated, so no
factor in the product formula for $\mathcal R_{K,r}(0)$ in
Theorem~\ref{thm:scalar-support} vanishes.  Every prime counted in
\eqref{eq:perratio} divides the nonzero integer
$|N_{\mathbf Q(\tau)/\mathbf Q}(c_0)|$: for such a prime, some
$\mathfrak q\mid p$ has $c_0\equiv0\pmod{\mathfrak q}$ by
Theorem~\ref{thm:scalar-support}, and $\mathfrak q$ has residue degree
one.  It remains to bound this norm.  Under the $j$-th real embedding,
$|\sigma_j(u_n)|=|\sin(2\pi jn/K)|/|\sin(2\pi j/K)|\le K/(4d_j)$ with
$d_j=\min(j,K/2-j)$, using $|\sin\pi x|\ge2\,\mathrm{dist}(x,\mathbf Z)$.
Since $\alpha_r=u_r$ and $\beta_r=-2u_{r-1}$, each of the $m^2$ factors of
$c_0$ has $j$-th embedding at most $2+K/d_j$ in absolute value, so
\[
\log|N(c_0)|
 \le m^2\sum_{j=1}^m\log(2+K/d_j)=O(K^3):
\]
the last sum is $O(K)$ by the elementary Stirling estimate applied to the
integers and half-integers $d_j$.  Summing $\log p$ over the distinct
counted primes, each dividing $|N(c_0)|$, gives \eqref{eq:perratio}.
Summing over $O(K)$ ratio classes recovers an $O(K^4)$ scale, so
\eqref{eq:perratio} is not a polynomial largest-prime theorem.

\section{Inverse reflection and the pair-sum square}

The ratio $r=-1$ has additional symmetry.  Here
\[
\alpha=-1,\qquad \beta=2\tau,
\qquad Q(T)=-P(2\tau-T).
\]
Set
\begin{equation}\label{eq:FEO}
F(X)=P(\tau+X)=E(X^2)+XO(X^2).
\end{equation}
The reflection reduction gives
\begin{equation}\label{eq:inverse-gcd}
I_{-1}(\mathfrak q)=2\deg\gcd(E,O\bmod\mathfrak q).
\end{equation}
Thus inverse-row intersections are automatically even.  A fixed-degree
half-size subresultant construction defines the exact split support integer
$\Xi_K$.

The scalar shift resultant admits a stronger universal factorization.

\begin{theorem}[Inverse pair-sum factorization]\label{thm:pairsum}
Let $x_1,\ldots,x_m$ be the roots of the monic polynomial $F$.  Define
\[
S(Z)=\prod_{i<j}(Z-x_i-x_j).
\]
Then $S\in\Z[\tau][Z]$ and
\begin{equation}\label{eq:pairsquare}
\Res_X(F(X),F(Z-X))=2^mF(Z/2)S(Z)^2.
\end{equation}
For every split embedding prime,
\begin{equation}\label{eq:inverse-order}
I_{-1}(\mathfrak q)=2\ord_{Z=0}(S(Z)\bmod\mathfrak q).
\end{equation}
Hence
\begin{equation}\label{eq:two-scalar}
I_{-1}(\mathfrak q)\ge3
\iff
S(0)\equiv S'(0)\equiv0\pmod{\mathfrak q}.
\end{equation}
\end{theorem}

\begin{proof}
The root-product formula gives
\[
\Res_X(F(X),F(Z-X))=\prod_{i,j}(Z-x_i-x_j).
\]
The diagonal ordered pairs give $2^mF(Z/2)$, and every unordered pair $i<j$
occurs twice.  The coefficients of $S$ are symmetric integral polynomials in
the roots.  At a split prime, $p\ne K$ is unramified and
$H_K\bmod\mathfrak q$ is separable -- its roots, the trace values $t_a$,
are pairwise distinct modulo $\mathfrak q$ -- so $\tau$ is a simple root
modulo $\mathfrak q$, i.e. $F(0)=P(\tau)=H'_K(\tau)\not\equiv0
\pmod{\mathfrak q}$ (and $\tau\not\equiv2\pmod{\mathfrak q}$).  Since $p$
is odd, the diagonal factor $2^mF(0)$ is a unit at $Z=0$.  Combining with the
general shift-resultant identity gives \eqref{eq:inverse-order};
\eqref{eq:two-scalar} follows because the left side is even.
\end{proof}

Let $S(Z)=s_0+s_1Z+\cdots$ and let the ideal $(s_0,s_1)$ contract to
$\Theta_K\Z$.  Then $\Theta_K$, $\Xi_K$, and the inverse factor of
$\Lambda_{K,-1}^{(3)}$ have the same split rational-prime support; their
valuations need not agree.

\section{The pair-sum polynomial as a translated half-degree resultant}

The square root $S$ itself can be obtained without constructing the
degree-$m^2$ shift resultant.  Translate
\begin{equation}\label{eq:translatedF}
F_Z(X)=F(X+Z/2)=E_Z(X^2)+XO_Z(X^2).
\end{equation}

\begin{theorem}[Translated parity-resultant identity]\label{thm:translated-res}
There exists a convention-dependent sign $\varepsilon_m\in\{\pm1\}$ such that
\begin{equation}\label{eq:translated-res}
S(Z)=\varepsilon_m\Res_Y(E_Z(Y),O_Z(Y)).
\end{equation}
Consequently
\begin{align}
S(0)&=\varepsilon_m\Res_Y(E,O),\label{eq:s0res}\\
S'(0)&=\varepsilon_m
\left.\frac{d}{dZ}\Res_Y(E_Z,O_Z)\right|_{Z=0}.\label{eq:s1res}
\end{align}
The first-order deformation is
\begin{equation}\label{eq:deformation}
\left.\partial_ZE_Z\right|_0=\frac{O+2YO'}{2},
\qquad
\left.\partial_ZO_Z\right|_0=E'.
\end{equation}
\end{theorem}

\begin{proof}
The roots of $F_Z$ are $y_i=x_i-Z/2$.  Apply the constant pair-sum identity
from Theorem~\ref{thm:pairsum} to the monic polynomial $F_Z$.  Its unordered
pair-sum product at zero is
\[
\prod_{i<j}(-y_i-y_j)
 =\prod_{i<j}(Z-x_i-x_j)=S(Z).
\]
It remains to identify this product with the parity resultant.  For a
monic polynomial $G$ with parity decomposition
$G(X)=\mathcal E(X^2)+X\mathcal O(X^2)$ and roots $y_i$, the constant
pair-sum coefficient satisfies $S_G(0)^2=\Res_Y(\mathcal E,\mathcal O)^2$.
Indeed, \eqref{eq:pairsquare} at $Z=0$ gives
$\Res_X(G(X),G(-X))=2^mG(0)S_G(0)^2$, while the root-product formula and
$G(-y_i)-G(y_i)=-2y_i\mathcal O(y_i^2)$ give
$\Res_X(G(X),G(-X))=\prod_iG(-y_i)
=2^m\mathcal E(0)\Res_Y(\mathcal E,\mathcal O)^2$:
here $\prod_iy_i=(-1)^m\mathcal E(0)$, and
$\prod_i\mathcal O(y_i^2)=\Res_X(G,\mathcal O(X^2))
=\Res_Y(\mathcal E,\mathcal O)^2$, evaluating at the roots $\pm\sqrt y$ of
$\mathcal O(X^2)$ and using that exactly one of $\mathcal E,\mathcal O$
is monic, in either parity case.  Since $G(0)=\mathcal E(0)\ne0$ in
characteristic zero, $S_G(0)^2=\Res_Y(\mathcal E,\mathcal O)^2$.  Apply
this to $G=F_Z$ over $\mathbf Q(\tau)(Z)$: its pair-sum product at zero
is $S(Z)$ by the display above, and its parity components are $E_Z,O_Z$
by \eqref{eq:translatedF}.  Hence
$S(Z)^2=\Res_Y(E_Z,O_Z)^2$ as polynomials in $Z$; both are nonzero in
the domain $\Z[\tau][Z]$, so they differ by a constant sign
$\varepsilon_m\in\{\pm1\}$, fixed by the monicity of $S$ once the
resultant convention is fixed.  Differentiating
$F(X+Z/2)$ at $Z=0$ gives \eqref{eq:deformation}.
\end{proof}

Thus inverse order-three support is the simultaneous divisibility problem for a
half-degree resultant and its translation derivative.  In exact computation,
both scalars can be obtained in one Sylvester determinant over dual numbers
$\varepsilon^2=0$, reducing the determinant dimension to $O(K)$ rather than
materializing a degree-$m^2$ polynomial.

Exact Smith contractions give, after removing factors $2K$,
\[
\Theta'_7=1,\quad
\Theta'_{11}=1,\quad
\Theta'_{13}=5,\quad
